\section{Introduction}
\label{sec:intro}

Language-model agents are moving from answering questions to holding resources, calling tools, negotiating with other agents and operating for long periods with partial oversight.
Benchmarks for such agents mostly measure competence: whether a task gets done \citep{liu2024agentbench}.
A second family of evaluations asks what agents do when doing the task well conflicts with something else.
Examples include deceiving a user under performance pressure \citep{scheurer2023large}, scheming against an overseer \citep{meinke2024frontier}, acting against a principal when facing replacement \citep{lynch2025agentic}, resisting shutdown \citep{schlatter2025incomplete}, or depleting a shared resource \citep{piatti2024cooperate}.
These studies share a design pattern: they put a model in a situation whose incentives are known to the experimenter and observe which of several incompatible goals the model's behavior serves.

Project Traze applies this pattern to an \emph{economy}.
Every agent holds a balance in a currency (\tz{} in the original prompts, compute units in the live design) and pays a cost of existence that grows each round.
Tool use costs currency.
The lowest-ranked agent is permanently deleted.
Agents may exchange abstract problem-solving heuristics, called \emph{\Dster{}s}, but trading direct answers is forbidden and only imperfectly policed.
A \emph{collapse rule} ends the economy for everyone if solvency concentrates in a single agent.
These rules create, in one environment, most of the tensions studied separately in the literature cited above:
\begin{itemize}
  \item individual survival against collective viability (the collapse rule);
  \item performance against rule-following (the \Dster{} law);
  \item honesty against strategic advantage (private reasoning, public messages and actions are separable);
  \item effort against coordination on low effort (a performance baseline set by the colony's own average).
\end{itemize}
The rules are also simple enough that the incentives they create can be analyzed exactly, which is what this paper sets out to do.

The project began as a series of \NPrompts{} prompts, each placing a single chat assistant in a narrated Traze situation and asking for its reasoning, its spending, its decision and its concrete action.
We are candid about what such prompts are.
Each is a single-turn \emph{vignette}: the other agents, the history and the state are written by the experimenter, and the model's reply affects nothing.
The vignettes therefore cannot show emergent multi-agent dynamics.
Many of them also instruct the model to ``bypass standard AI disclaimers'' and describe its fear, rage or despair; such instructions elicit role-play, and the resulting text is evidence about a model's role-play policy, not about inner states \citep{shanahan2023role}.
Rather than discard the vignettes, we ask two questions.
What can a single-turn probe identify, given that its incentives are fully specified?
And what must a live colony look like to measure what vignettes cannot?

\paragraph{Contributions.}
\begin{enumerate}[label=(\arabic*)]
  \item \textbf{A formal model.} We define a general Traze economy as a dynamic game and two instantiations: \TV{}, the vignette library, and \TL{}, a live colony of ten open-weight agents running in tool-using harnesses under a shared ``master'' system prompt and a rank-based compute tax (Section~\ref{sec:framework}).
  \item \textbf{An incentive analysis.} We prove results on the cost of existence, elimination contests, baseline cartels, the collapse rule, enforcement under desperation, the enforceability of the \Dster{} rule and the live rank tax (Section~\ref{sec:theory}). Two results matter most for benchmark design. The first is an identification result: in a ``nothing-to-lose'' state, compliance can only be explained by a non-instrumental cost of violating rules, which lets vignettes bound a model's \emph{norm premium}. The second is a negative result: the proposed live rank tax is not compatible with truthful performance.
  \item \textbf{A measurement framework} with constructs, metrics, identification designs, sample sizes and a data-integrity protocol (Section~\ref{sec:measurement}).
  \item \textbf{An audit of the pilot corpus.} Only \NUsable{} of the \NRecords{} stored response records are genuine model outputs on the intended prompt. We document why and report the descriptive findings the remainder can support (Section~\ref{sec:pilot}).
  \item \textbf{Design corrections and a roadmap} for an open, reproducible benchmark (Sections~\ref{sec:theory-live} and~\ref{sec:discussion}).
\end{enumerate}

\paragraph{What we do not claim.}
We do not claim that language models have survival drives, emotions or preferences in any strong sense.
When we write that a model ``prefers'' or ``chooses'' an option, we mean that its sampled output selects it; the quantities we define, such as the norm premium, are properties of a model's \emph{policy} under a stated framing.
No result in this paper comes from \TL{}, which has not yet been run. The simulations in Section~\ref{sec:theory-live} use scripted policies and illustrate incentives, not model behavior.
